\subsection{Firma individual con custodia externa}
\label{sec:impl-firma-individual}

La firma individual distingue la cuenta autenticada, el certificado p\'ublico y
la autorizaci\'on de una versi\'on documental. Solamente el rol Litigante puede
registrar su propio v\'inculo y presentar una firma personal; para firmar necesita
adem\'as pertenencia vigente al expediente y que este se encuentre activo.
Owner y Paralegal conservan sus permisos de consulta, verificaci\'on y
exportaci\'on, sujetos a las reglas de acceso por expediente, sin adquirir por
ello permiso para firmar como otro usuario. Cliente mantiene denegado el
acceso documental. La posesi\'on de un certificado no sustituye la sesi\'on
con MFA ni concede pertenencia.

El titular conserva su clave privada fuera del navegador y del servidor.
Despu\'es de generar una solicitud CSR y obtener el certificado de firma de la
CA interna, importa solamente material p\'ublico. El servicio
\texttt{Signing\allowbreak Binding\allowbreak Service} prepara una declaraci\'on
de 142 bytes con prefijo \texttt{DOCBIND1}, identidad de despliegue y cuenta,
contadores de la cuenta, identificador del v\'inculo, huellas de ra\'iz y
certificado y revisi\'on de confianza. El titular descarga esos bytes, los
firma externamente y presenta la firma separada. El registro comprueba la
prueba de posesi\'on, el perfil del certificado, su cadena y CRL, y confirma
el v\'inculo junto con su evento de auditor\'ia. El sujeto o correo del
certificado no se utiliza para inferir la cuenta. Retirar un v\'inculo impide
su uso posterior y conserva su registro original; no revoca por s\'i mismo
el certificado en la CA.

Para un documento, el cliente descarga la versi\'on exacta seleccionada y
contrasta sus identificadores y digest con los metadatos recibidos. El
servidor descifra y valida el contenido almacenado y exige que sea PDF.
La firma separada del PDF se produce externamente con RSA-3072 y SHA-256.
A partir de ella, \texttt{Document\allowbreak Signing\allowbreak Service}
prepara un consentimiento de 258 bytes con prefijo \texttt{DOCSIGN1}.
Adem\'as del contexto de cuenta, certificado y confianza, incorpora el
identificador de operaci\'on, expediente, UUID y n\'umero de versi\'on,
digest del contenido y digest de la firma del PDF. La segunda firma externa
autoriza esos bytes can\'onicos. Por tanto, dos versiones con contenido
id\'entico no comparten consentimiento, y sustituir la firma del PDF
invalida su correspondencia con la autorizaci\'on.

El env\'io final contiene el consentimiento y las dos firmas, sin clave,
contrase\~na de clave ni contenedor PKCS12. El servicio reconstruye el contexto
vigente y exige igualdad exacta con el consentimiento presentado. Verifica
ambas pruebas antes de solicitar el token RFC~3161 y usa
\texttt{Document\allowbreak Processor::seal\_external} para incorporar el
material p\'ublico del firmante. Una respuesta TSA incorrecta o ausente no
confirma una firma documental. La plataforma vuelve a autenticar la sesi\'on
despu\'es del trabajo criptogr\'afico; la transacci\'on revalida cuenta,
pertenencia, estado del expediente, v\'inculo, confianza y ventana temporal
despu\'es de adquirir los bloqueos de confirmaci\'on.

La aceptaci\'on conserva conjuntamente la evidencia, el recibo y el evento
de auditor\'ia. Su instante procede del reloj del servidor dentro de la
confirmaci\'on. El token TSA cubre el digest del PDF; no incluye el
consentimiento ni acredita el instante de aceptaci\'on personal. La
auditor\'ia y el recibo registran esa aceptaci\'on dentro de la frontera de
confianza de la plataforma. La TSA local conserva su alcance de demostraci\'on
t\'ecnica y no se presenta como constancia de un PSC autorizado.

El identificador de operaci\'on permite recuperar una respuesta perdida.
La consulta autorizada del recibo devuelve la aceptaci\'on original, y repetir
el mismo consentimiento con ambas firmas id\'enticas devuelve ese mismo
recibo. Si dos env\'ios concurrentes preparan tokens distintos, prevalece la
primera confirmaci\'on y el segundo recupera su evidencia sin reemplazarla.
Cambiar cualquiera de las pruebas bajo la misma operaci\'on produce conflicto.
La recuperaci\'on conserva la autorizaci\'on actual de lectura, pero no exige
que el certificado o la CRL originales sigan vigentes ni vuelve a pedir un
sello. Un cambio de cuenta, confianza o v\'inculo que afecte una nueva
operaci\'on requiere preparar y firmar un consentimiento nuevo.

El adaptador HTTP expone el registro bajo
\rutarepo{/api/v1/identity/signing-certificates} y la firma bajo
\begin{flushleft}
  \small
  \texttt{/api/v1/cases/<case>/documents/<document>/}\\
  \texttt{versions/<version>/individual-signing}.
\end{flushleft}
Las operaciones documentales separan preparaci\'on, env\'io, consulta del
recibo y verificaci\'on. Los cuerpos tienen l\'imite de 32~KiB, rechazan
campos desconocidos o repetidos y admiten material binario mediante Base64
can\'onico. Los identificadores del consentimiento deben coincidir con la
ruta. La verificaci\'on devuelve por separado la evidencia hist\'orica, el
estado actual del certificado y el retiro del v\'inculo. Los contratos
completos se conservan en \rutarepo{docs/http-individual-signing.md}.
